\documentclass[11pt]{article}

\usepackage[margin=1in]{geometry}
\usepackage{amsmath,amssymb}
\usepackage{graphicx}
\usepackage{hyperref}
\usepackage{booktabs}
\usepackage{enumitem}
\usepackage{xcolor}

\hypersetup{
    colorlinks=true,
    linkcolor=blue,
    citecolor=blue,
    urlcolor=blue
}

\newcommand{\Ec}{E_{\mathrm{c}}}
\newcommand{\Ehad}{E_{\mathrm{had}}}
\newcommand{\Emu}{E_{\mu}}
\newcommand{\EHNL}{E_{N}}
\newcommand{\fEM}{f_{\mathrm{EM}}}
\newcommand{\CCH}{C_{\mathrm{Ch}}}
\newcommand{\Adet}{A_{\mathrm{det}}}
\newcommand{\Rdet}{R_{\mathrm{det}}}
\newcommand{\shad}{\sigma_{\mathrm{had}}}
\newcommand{\thetaC}{\theta_{\mathrm{C}}}

\title{Analytical Hadronic Shower Cherenkov Model\\
for Balloon HNL Detection}
\author{N.\ Kamp\thanks{Model implementation assisted by Claude (Anthropic).}}
\date{\today}

\begin{document}
\maketitle

\begin{abstract}
We describe an analytical model for atmospheric Cherenkov light from
hadronic showers, developed for a balloon-based search for heavy neutral
lepton (HNL) decays at a neutrino factory.  The model combines three
ingredients from the imaging atmospheric Cherenkov telescope (IACT)
literature: (i)~the calorimetric track-length integral, which gives a
total Cherenkov photon yield proportional to shower energy and
independent of altitude; (ii)~an electromagnetic-fraction scaling for
hadronic showers via the Gaisser parameterization; and (iii)~a Gaussian
geometric acceptance that exploits the fact that the hadronic angular
spread ($\shad \sim 10^\circ$) greatly exceeds the Cherenkov angle
($\thetaC \sim 1.4^\circ$).  The resulting closed-form expression
replaces a Monte Carlo sub-track simulation, yielding identical physics
with substantially reduced computation.
\end{abstract}

\tableofcontents

%======================================================================
\section{Introduction}
\label{sec:intro}
%======================================================================

A neutrino factory directs $\sim 10^{22}$ muon decays over a ten-year
lifetime into a beam dump embedded in the Earth.  Deep-inelastic
scattering (DIS) of the muon beam, $\mu + N \to N_4 + X$, can produce
heavy neutral leptons (HNLs) with masses up to tens of GeV.  The HNLs
traverse the rock unimpeded and decay in the atmosphere; the charged
products emit Cherenkov light detectable by an elevated camera.

For HNL masses above $\sim 1$~GeV, the dominant decay channel is the
charged-current mode
\begin{equation}
  N_4 \to \mu + \text{hadrons}\,,
  \label{eq:hnl_decay}
\end{equation}
producing a muon and a hadronic system.  Previous versions of the
simulation modeled only the muon-track Cherenkov emission and either
ignored the hadronic component or attempted to track individual decay
channels.  This note documents an analytical model that captures the
hadronic Cherenkov contribution without per-particle ray-tracing.

The same formalism applies to the neutrino background: CC DIS events
$\nu_\mu + N \to \mu + X$ in the atmosphere also produce a muon plus a
hadronic shower, and both components contribute Cherenkov light.

%======================================================================
\section{Cherenkov Yield from Electromagnetic Showers}
\label{sec:em_yield}
%======================================================================

\subsection{Frank--Tamm formula}

A relativistic charged particle ($\beta \approx 1$) traversing a medium
with refractive index $n$ emits Cherenkov photons at a rate given by the
Frank--Tamm formula:
\begin{equation}
  \frac{dN}{d\ell} \;=\;
  2\pi\alpha\,\sin^2\!\thetaC
  \int_{\lambda_{\min}}^{\lambda_{\max}}
  \frac{d\lambda}{\lambda^2}
  \;=\;
  2\pi\alpha\,\sin^2\!\thetaC
  \left(\frac{1}{\lambda_{\min}} - \frac{1}{\lambda_{\max}}\right),
  \label{eq:frank_tamm}
\end{equation}
where $\alpha = 1/137$ is the fine-structure constant and
$\thetaC = \arccos(1/n)$ is the Cherenkov angle.  For air at sea level
($n = 1.0003$), with $\lambda_{\min} = 300$~nm and
$\lambda_{\max} = 1000$~nm:
\begin{equation}
  \sin^2\!\thetaC \;\approx\; 2(n-1) \;=\; 6.0 \times 10^{-4},
  \qquad
  \left.\frac{dN}{d\ell}\right|_{\text{sea}}
  \;\approx\; 64\;\text{photons\,m}^{-1}.
  \label{eq:dndl_sea}
\end{equation}

\subsection{Track-length integral}

An electromagnetic (EM) shower of primary energy $E$ develops through
bremsstrahlung and pair production.  At shower maximum, the number of
charged particles above the critical energy $\Ec$ is approximately
$E/\Ec$, and each particle travels roughly one radiation length $X_0$
before losing energy.  The total Cherenkov-producing track length of all
shower particles is therefore~\cite{Nerling2006,IceCube2013}
\begin{equation}
  T_{\mathrm{total}}(E,z) \;=\;
  \frac{E}{\Ec}\;\frac{X_0}{\rho(z)}\,,
  \label{eq:track_length}
\end{equation}
where $X_0 = 36.62\;\text{g\,cm}^{-2}$ is the radiation length of air,
$\Ec = 81$~MeV is the critical energy, and $\rho(z)$ is the air
density at altitude $z$.  The total number of Cherenkov photons emitted
by the shower is
\begin{equation}
  N_{\gamma}^{\mathrm{EM}}(E) \;=\;
  T_{\mathrm{total}} \times \frac{dN}{d\ell}(z)
  \;=\;
  \frac{E}{\Ec}\;\frac{X_0}{\rho(z)}
  \times \frac{dN}{d\ell}(z)\,.
  \label{eq:Nph_em}
\end{equation}

\subsection{Altitude independence}

The Cherenkov yield per meter scales with density through the refractive
index:
\begin{equation}
  \frac{dN}{d\ell}(z) \;\propto\; \sin^2\!\thetaC(z) \;\propto\;
  \bigl[n(z) - 1\bigr] \;\propto\; \rho(z)\,.
\end{equation}
Since $T_{\mathrm{total}} \propto 1/\rho(z)$ from
Eq.~\eqref{eq:track_length}, the product is independent of altitude:
\begin{equation}
  \boxed{
  N_{\gamma}^{\mathrm{EM}}(E)
  \;=\;
  \frac{E}{\Ec}\;\CCH\,,
  \qquad
  \CCH \;\equiv\;
  \frac{X_0}{\rho_0}
  \times
  \left.\frac{dN}{d\ell}\right|_{\text{sea}}\,.
  }
  \label{eq:C_CH}
\end{equation}
Numerically, with $\rho_0 = 1.225 \times 10^{-3}\;\text{g\,cm}^{-3}$:
\begin{equation}
  \frac{X_0}{\rho_0} = 299\;\text{m}\,,
  \qquad
  \CCH \approx 299 \times 64 \approx 19{,}200\;\text{photons}\;\text{per}\;(E/\Ec).
  \label{eq:C_CH_value}
\end{equation}
Equivalently, $N_{\gamma}^{\mathrm{EM}} \approx 2.4 \times 10^5 \times
E\,[\text{GeV}]$ photons are emitted in total by the shower, regardless
of altitude.  This result is validated by the IceCube Geant4
study~\cite{IceCube2013}, which found a linear relationship
$\hat\ell = \alpha\,E$ to precision $10^{-5}$ for EM cascades.

%======================================================================
\section{Extension to Hadronic Showers}
\label{sec:hadronic}
%======================================================================

\subsection{EM-fraction scaling}

A hadronic shower cascades through successive nuclear interactions.  At
each interaction vertex, roughly one-third of the produced pions are
neutral ($\pi^0 \to \gamma\gamma$), transferring energy irreversibly
into the electromagnetic channel.  After $k$ generations, the
electromagnetic fraction is $1 - (2/3)^k$.  This is well described by
the Gaisser parameterization~\cite{Gaisser1990,Wigmans2000}:
\begin{equation}
  \fEM(E) \;=\;
  1 - \left(\frac{E}{1\;\text{GeV}}\right)^{-0.14}.
  \label{eq:f_EM}
\end{equation}
The exponent $-0.14 \approx \ln(2/3)/\ln(E_0/E)$ encodes the
multiplicity growth of the hadronic cascade.  At 5~GeV,
$\fEM \approx 0.20$; at 50~GeV, $\fEM \approx 0.43$.

The total Cherenkov photon yield from a hadronic shower of energy
$\Ehad$ is then
\begin{equation}
  N_{\gamma}^{\mathrm{had}}(\Ehad)
  \;=\;
  \fEM(\Ehad)\;\frac{\Ehad}{\Ec}\;\CCH\,.
  \label{eq:Nph_had}
\end{equation}

%======================================================================
\section{Geometric Acceptance}
\label{sec:geometry}
%======================================================================

\subsection{Angular spread of hadronic showers}

The charged particles in a hadronic shower acquire transverse momentum
from nuclear interactions.  The characteristic transverse momentum is
$p_T \sim 350$~MeV, giving an angular spread relative to the shower
axis of order
\begin{equation}
  \shad \;\sim\; \frac{p_T}{\Ehad / N_{\mathrm{ch}}}
  \;\sim\; 10^\circ
\end{equation}
for GeV-scale showers, where $N_{\mathrm{ch}}$ is the charged-particle
multiplicity.  This is much larger than the Cherenkov angle in air:
\begin{equation}
  \thetaC \;=\; \arccos(1/n) \;\approx\;
  \sqrt{2(n-1)} \;\approx\; 1.4^\circ
  \quad\text{(sea level)}.
\end{equation}

\subsection{Smearing of the Cherenkov ring}

Each charged particle in the shower emits Cherenkov photons in a thin
cone of half-angle $\thetaC$ around its direction of motion.  When the
particle directions are distributed as a Gaussian with width $\shad$
around the shower axis, and $\shad \gg \thetaC$, the Cherenkov ring
from each particle is narrow compared to the spread of emission
directions.  Integrating over all particles and their Cherenkov
azimuths, the resulting angular distribution of photons relative to the
shower axis is well approximated by the particle distribution itself:
\begin{equation}
  P(\alpha) \;\approx\;
  \frac{1}{2\pi\,\shad^2}\,
  \exp\!\left(-\frac{\alpha^2}{2\,\shad^2}\right),
  \label{eq:P_alpha}
\end{equation}
where $\alpha$ is the angle between the shower axis $\hat{d}$ and the
direction from the shower to the observation point.  This is normalized
on the sphere in the small-angle limit:
$\int_0^\infty 2\pi\alpha\, P(\alpha)\,d\alpha = 1$.

\subsection{Detected photon count}

A detector of area $\Adet = \pi\Rdet^2$ at distance $d$ from the
shower subtends solid angle $\Omega_{\mathrm{det}} = \Adet / d^2$.  If
the detector is at angle $\alpha$ from the shower axis, the fraction of
photons intercepted is
\begin{equation}
  \Phi(\alpha, d)
  \;=\;
  P(\alpha) \times \frac{\Adet}{d^2}
  \;=\;
  \frac{\Adet}{2\pi\,\shad^2\,d^2}\,
  \exp\!\left(-\frac{\alpha^2}{2\,\shad^2}\right).
  \label{eq:geom}
\end{equation}

%======================================================================
\section{Complete Analytical Formula}
\label{sec:formula}
%======================================================================

Combining the hadronic yield (Eq.~\ref{eq:Nph_had}), the geometric
acceptance (Eq.~\ref{eq:geom}), and atmospheric transmission $\mathcal{T}(z,\theta_z)$
accounting for Rayleigh and Mie scattering along the slant path from
altitude $z$ to the detector, the number of detected Cherenkov photons
from a hadronic shower is
\begin{equation}
  \boxed{
  N_{\mathrm{det}}^{\mathrm{had}}
  \;=\;
  \fEM(\Ehad)\;\frac{\Ehad}{\Ec}\;\CCH
  \;\times\;
  \frac{\Adet}{2\pi\,\shad^2\,d^2}\;
  e^{-\alpha^2/2\shad^2}
  \;\times\;
  \mathcal{T}(z, \theta_z)\,.
  }
  \label{eq:master}
\end{equation}
The key features of this formula are:
\begin{itemize}[nosep]
  \item \textbf{Linear in energy}: $N_{\mathrm{det}} \propto \Ehad$
    (modulo the slowly-varying $\fEM$).
  \item \textbf{Altitude-independent yield}: the $\CCH$ constant
    absorbs all density dependence. The only altitude-dependent factor
    is the atmospheric transmission $\mathcal{T}$.
  \item \textbf{Analytical}: no ray-tracing or Monte Carlo sampling
    is required for the hadronic shower component.
  \item \textbf{Gaussian angular profile}: the detector response falls
    as $e^{-\alpha^2/2\shad^2}$, favoring on-axis geometries.
\end{itemize}

\noindent
\textbf{Numerical constants used in the implementation:}

\begin{center}
\begin{tabular}{lll}
\toprule
Symbol & Value & Description \\
\midrule
$\Ec$ & 80 MeV & Critical energy in air \\
$X_0$ & 36.62 g\,cm$^{-2}$ & Radiation length of air \\
$\rho_0$ & $1.225 \times 10^{-3}$ g\,cm$^{-3}$ & Sea-level air density \\
$n$ & 1.0003 & Refractive index of air (sea level) \\
$\lambda_{\min},\,\lambda_{\max}$ & 300, 1000 nm & Cherenkov wavelength band \\
$dN/d\ell|_{\text{sea}}$ & 64 m$^{-1}$ & Cherenkov yield at sea level \\
$\CCH$ & 19,200 & Photons per $(E/\Ec)$ \\
$\shad$ & $10^\circ$ & Hadronic angular spread \\
$\Rdet$ & 2 m & Detector radius \\
\bottomrule
\end{tabular}
\end{center}

%======================================================================
\section{Application to HNL Signal}
\label{sec:signal}
%======================================================================

The HNL decay $N_4 \to \mu + \text{hadrons}$ produces a muon of energy
$\Emu$ and a hadronic system of energy $\Ehad = \EHNL - \Emu$.  Both
are available from the pre-computed Monte Carlo data tables, which
contain the full four-momenta $(E_N, \vec{p}_N)$ and
$(\Emu, \vec{p}_\mu)$ for each simulated decay.

The hadronic shower direction is determined by momentum conservation:
\begin{equation}
  \vec{p}_{\mathrm{had}} \;=\;
  \EHNL\,\hat{d}_N - \Emu\,\hat{d}_\mu\,,
  \qquad
  \hat{d}_{\mathrm{had}} \;=\;
  \frac{\vec{p}_{\mathrm{had}}}{|\vec{p}_{\mathrm{had}}|}\,.
  \label{eq:had_dir}
\end{equation}
The total detected photon count for each MC event is then
\begin{equation}
  N_{\mathrm{det}}^{\mathrm{total}}
  \;=\;
  N_{\mathrm{det}}^{\mu}
  \;+\;
  N_{\mathrm{det}}^{\mathrm{had}}\,,
\end{equation}
where $N_{\mathrm{det}}^{\mu}$ is computed from the muon track via
standard Cherenkov ray-tracing (the muon has no angular spread, so the
narrow-cone geometry must be preserved), and
$N_{\mathrm{det}}^{\mathrm{had}}$ is given by
Eq.~\eqref{eq:master}.

%======================================================================
\section{Application to Neutrino Background}
\label{sec:background}
%======================================================================

The beam-induced neutrino background arises from the same DIS process
that produces HNLs, but with zero mass:
$\mu + N \to \nu_\mu + X$.  These neutrinos travel upward and interact
in the atmosphere via charged-current DIS:
\begin{equation}
  \nu_\mu + N \to \mu + X\,.
\end{equation}
The inelasticity $y$ is sampled from
$d\sigma/dy \propto 1 + (1-y)^2$, giving
\begin{equation}
  \Emu^{\mathrm{out}} = (1-y)\,E_\nu\,,
  \qquad
  \Ehad = y\,E_\nu\,.
\end{equation}
The hadronic shower direction is taken as the neutrino direction
(collinear approximation), and the muon direction is scattered by a
characteristic angle $\sim m_\mu / \Emu^{\mathrm{out}}$ from the
neutrino.  The detected photon count is again
\begin{equation}
  N_{\mathrm{det}}^{\mathrm{bg}}
  \;=\;
  N_{\mathrm{det}}^{\mu,\mathrm{bg}}
  \;+\;
  N_{\mathrm{det}}^{\mathrm{had,bg}}\,,
\end{equation}
with $N_{\mathrm{det}}^{\mathrm{had,bg}}$ from Eq.~\eqref{eq:master}
using $\Ehad = y\,E_\nu$ and the neutrino direction as the shower axis.

Each background event carries the CC interaction weight
$w = \sigma_{\mathrm{CC}}(E_\nu) \times
\mathcal{N}_{\mathrm{col}}(z_{\max})$,
representing the probability that the neutrino interacts in the
atmospheric column below the detector.

%======================================================================
\section{Numerical Validation}
\label{sec:validation}
%======================================================================

\subsection{Yield constant}

The computed value of $\CCH$ from the code is
\begin{equation}
  \CCH = \frac{X_0}{\rho_0}
  \times \left.\frac{dN}{d\ell}\right|_{\text{sea}}
  = 298.9\;\text{m} \times 64.2\;\text{m}^{-1}
  = 19{,}181\;\text{photons}\;\text{per}\;(E/\Ec),
\end{equation}
consistent with the estimate of $\sim 19{,}200$.

\subsection{Example: 5~GeV hadronic shower}

Consider a 5~GeV hadronic shower at altitude $z = 5$~km, with a
detector directly above at $z = 10$~km (distance $d = 5$~km, on-axis
$\alpha = 0$):
\begin{align}
  \fEM(5) &= 1 - 5^{-0.14} = 0.202\,, \\
  N_{\gamma}^{\mathrm{had}} &= 0.202 \times \frac{5}{0.080}
    \times 19{,}181 = 242{,}000\;\text{photons}\,, \\
  \Phi(0^\circ) &= \frac{\pi \times 2^2}{2\pi \times
    (10\pi/180)^2 \times (5000)^2}
    = 2.6 \times 10^{-6}\,, \\
  N_{\mathrm{det}} &\approx 242{,}000 \times 2.6 \times 10^{-6}
    \times \mathcal{T}
    \approx 0.6\,\mathcal{T}\;\text{photons}\,.
\end{align}
With $\mathcal{T} \sim 0.8$ at 5~km altitude and moderate zenith angle,
this gives $N_{\mathrm{det}} \approx 0.5$ photons.  While sub-threshold
on its own, this adds to the muon-track contribution and can push
borderline events above the detection threshold.

\subsection{Angular dependence}

The Gaussian suppression was verified against the analytical prediction:

\begin{center}
\begin{tabular}{ccc}
\toprule
$\alpha$ & $N_{\mathrm{det}}/N_{\mathrm{det}}(0)$ (code) &
$\exp(-\alpha^2/2\shad^2)$ \\
\midrule
$0^\circ$  & 1.000 & 1.000 \\
$5^\circ$  & 0.881 & 0.882 \\
$15^\circ$ & 0.320 & 0.325 \\
$30^\circ$ & 0.010 & 0.011 \\
\bottomrule
\end{tabular}
\end{center}

\noindent
The small deviations at large $\alpha$ arise from the
altitude-dependent atmospheric transmission varying with zenith angle.

%======================================================================
\section{Summary}
\label{sec:summary}
%======================================================================

The analytical hadronic shower Cherenkov model provides a closed-form
expression (Eq.~\ref{eq:master}) for the detected photon count from
hadronic showers in the atmosphere.  The key physical ingredients are:

\begin{enumerate}[nosep]
  \item The \textbf{calorimetric track-length integral}: total
    Cherenkov yield $\propto E$, altitude-independent.
  \item The \textbf{EM-fraction scaling}: hadronic showers produce
    Cherenkov light proportional to their EM content,
    $\fEM = 1 - E^{-0.14}$.
  \item The \textbf{Gaussian geometric acceptance}: valid when
    $\shad \gg \thetaC$, replacing per-particle ray-tracing.
\end{enumerate}

\noindent
This model is applied identically to both HNL signal decays
($N_4 \to \mu + \text{hadrons}$) and neutrino DIS background events
($\nu_\mu + N \to \mu + X$), ensuring consistent treatment of the
hadronic Cherenkov contribution throughout the sensitivity analysis.

\begin{thebibliography}{10}

\bibitem{Nerling2006}
F.~Nerling, J.~Bl\"umer, R.~Engel, and M.~Risse,
``Universality of electron distributions in high-energy air showers ---
Description of Cherenkov light production,''
\textit{Astropart.\ Phys.}\ \textbf{24}, 421 (2006),
\href{https://arxiv.org/abs/astro-ph/0506729}{arXiv:astro-ph/0506729}.

\bibitem{IceCube2013}
IceCube Collaboration,
``Calculation of the Cherenkov light yield from electromagnetic
cascades in ice with Geant4,''
\textit{Astropart.\ Phys.}\ \textbf{44}, 102 (2013),
\href{https://arxiv.org/abs/1210.5140}{arXiv:1210.5140}.

\bibitem{Yuan2025}
T.~Yuan \textit{et al.},
``Probabilistic modeling of Cherenkov emission from particle showers,''
(2025),
\href{https://arxiv.org/abs/2601.00944}{arXiv:2601.00944}.

\bibitem{Gaisser1990}
T.~K.~Gaisser,
\textit{Cosmic Rays and Particle Physics},
Cambridge University Press (1990).

\bibitem{Wigmans2000}
R.~Wigmans,
\textit{Calorimetry: Energy Measurement in Particle Physics},
Oxford University Press (2000).

\bibitem{FrankTamm1937}
I.~Frank and I.~Tamm,
``Coherent visible radiation of fast electrons passing through matter,''
\textit{Dokl.\ Akad.\ Nauk SSSR}\ \textbf{14}, 109 (1937).

\end{thebibliography}

\end{document}
